\documentclass{../nndlcourse}

\begin{document}

\title{机器学习概述}
\author{数据科学与大数据技术专业}
\date{讲义1}
\maketitle

\section{导读}

除本导读外, 本讲共有 3 个主体节. 它们从“什么是机器学习”出发, 说明机器学习要解决哪些任务, 一个学习系统由哪些基本要素组成, 以及为什么模型不能只追求训练集表现.

\begin{enumerate}
  \item \textbf{机器学习的核心概念与任务分类}: 说明为什么需要机器学习, 它与人工智能和统计学的关系, 以及监督学习、无监督学习和强化学习的基本区别.
  \item \textbf{机器学习的四个基本要素}: 说明一个学习系统通常由数据、模型、学习准则和优化算法共同构成.
  \item \textbf{泛化能力与正则化}: 说明为什么训练集表现不等于真实能力, 以及正则化和验证集如何帮助控制过拟合.
\end{enumerate}

\textbf{如果细节都忘了}: 最应该记住的是, 机器学习不是把规则手写出来, 而是用数据选择一个能在新样本上工作的函数. 因此, 学习时必须同时问四件事: 数据是什么, 模型能表达什么, 损失函数奖励什么, 优化算法能不能找到好参数.

\textbf{问题思考}:

\begin{itemize}
  \item 如果一个系统在训练数据上表现很好, 为什么还不能直接说明它真正学会了任务规律?
\end{itemize}

\section{机器学习的核心概念与任务分类}

\textbf{本节导读}:
\begin{itemize}
  \item \textbf{本节目的}: 建立机器学习的基本定义和任务版图, 让后续算法都能放回同一个框架中理解.
  \item \textbf{为什么学}: 如果不知道任务类型和学习目标, 后面讨论模型、损失和优化时就很容易只记公式, 不知道它们在解决什么问题.
  \item \textbf{核心内容}: 机器学习的核心是从数据中学习规律, 并把这种规律推广到未见样本; 不同任务的区别主要体现在数据是否有标签、输出是什么以及反馈如何获得.
  \item \textbf{最小例子}: 垃圾邮件过滤很难靠人工列出所有规则, 但可以用大量“垃圾/非垃圾”邮件训练分类器, 让模型从词语、发件人和链接等特征中学习判别规律.
  \item \textbf{最该记住}: 机器学习的目标不是记住训练样本, 而是学到能迁移到新样本的输入输出关系或数据结构.
\end{itemize}

\begin{enumerate}
  \item \textbf{为什么需要机器学习}
    \begin{itemize}
      \item 在许多任务中, 人工设计精确规则十分困难, 例如视觉识别、语音处理和自然语言理解等问题往往具有高维、复杂和非结构化的特点.
      \item 真实环境通常处于持续变化之中, 系统需要能够根据新数据不断调整, 例如垃圾邮件过滤、推荐系统和风险控制.
      \item 某些问题涉及大规模模式发现与复杂决策, 仅依赖人工编写规则往往难以达到足够的性能上限.
      \item 在排序、资源分配和个性化服务等任务中, 通过数据驱动学习还可以更系统地兼顾效果、公平性与适应性.
    \end{itemize}
  \item \textbf{机器学习与相关领域的关系}
    \begin{itemize}
      \item 机器学习与统计学都关注从数据中提取规律, 并共同依赖概率论、线性代数与优化方法等数学基础.
      \item 二者的侧重点并不完全相同: 统计学通常更强调推断、解释与结论的可靠性, 机器学习通常更强调预测性能、计算效率与模型的可扩展性.
      \item 机器学习是人工智能的重要实现路径之一, 但人工智能并不等同于机器学习.
      \item 基于规则的系统、符号推理和搜索方法同样属于人工智能的重要组成部分, 它们并不一定依赖数据驱动的学习过程.
      \item 机器学习也常从人类学习和生物系统中获得启发, 但其目标是解决特定任务, 因而不必在机制上完全模拟人类认知过程.
    \end{itemize}
  \item \textbf{机器学习的基本定义}
    \begin{itemize}
      \item 机器学习是研究如何利用数据自动学习规律, 从而对未知样本进行预测、分类、决策或生成的理论、方法与技术.
      \item 从形式上看, 所学习的对象通常可以表示为函数关系、条件概率分布、决策规则或策略函数.
    \end{itemize}
  \item \textbf{主要任务类型}
    \begin{itemize}
      \item 监督学习: 利用带标签样本学习输入到输出的映射关系, 并将该映射推广到新样本.
        \begin{itemize}
          \item 回归任务: 预测连续型输出, 例如房价估计、销量预测和温度预测.
          \item 分类任务: 预测离散类别标签, 例如手写数字识别、邮件分类和疾病辅助诊断.
        \end{itemize}
      \item 无监督学习: 在无标签数据上发现潜在结构、表示或分布特征.
        \begin{itemize}
          \item 聚类: 按相似性将样本划分为若干组, 例如客户分群和图像分割.
          \item 降维: 将高维数据映射到低维表示, 便于可视化、压缩与后续建模, 例如主成分分析.
        \end{itemize}
      \item 强化学习: 智能体通过与环境交互获得反馈, 学习能够最大化长期累计回报的策略.
    \end{itemize}
  \item \textbf{简要发展脉络}
    \begin{itemize}
      \item 从人工智能的整体历史看, 机器学习并不是一开始就占据中心位置. 1956 年达特茅斯会议提出 ``Artificial Intelligence'' 这一名称时, 研究者更相信显式规则、逻辑推理和搜索方法能够模拟人的智能. 这种路线通常称为符号主义: 只要把知识写成足够多的 ``if-then'' 规则, 机器似乎就可以进行推理.
      \item 第一波热潮很快遇到边界. 一方面, 现实世界中的知识充满例外, 很难全部写成规则; 另一方面, 早期神经网络中的单层感知器只能表达线性可分问题, 连 XOR 这样的简单非线性分类也无法解决. 当过高承诺无法兑现时, 研究经费和社会期待迅速收缩, 形成了 20 世纪 70 年代的第一次 AI 寒冬.
      \item 20 世纪 80 年代, 专家系统让人工智能再次升温. 它把医生、工程师等领域专家的经验整理成规则库, 在诊断、配置和决策支持等垂直任务中取得过很亮眼的结果. 但专家系统的知识获取和维护成本极高, 新情况一旦超出规则库就难以处理. 这种 ``知识瓶颈'' 导致第二次 AI 寒冬, 也让研究者更深刻地意识到: 许多智能行为不能只靠人工编写规则, 还必须让系统从数据中学习.
      \item 在寒冬时期, 统计学习、模式识别和神经网络研究并没有完全停止, 只是常常以更低调的名字继续发展. 支持向量机、概率图模型、集成学习和反向传播等方法, 逐步把问题转向可计算、可优化、可评估的学习框架. 这也是本课程从数据、模型、学习准则和优化算法讲起的原因: 现代机器学习首先要把 ``智能'' 拆成可以训练和验证的具体任务.
      \item 进入 21 世纪后, 互联网积累了海量数据, GPU 等硬件提供了强大算力, 优化方法和开源工程生态也日趋成熟. 深度学习因此在图像识别、语音识别、机器翻译、博弈和大语言模型等任务中快速突破, 从 2012 年 AlexNet 到 AlphaGo, 再到 ChatGPT, 机器学习逐渐成为现代人工智能最重要的技术基础之一.
      \item 这段历史给我们的一个重要提醒是: 人工智能的发展并不是线性上升的故事, 而是 ``突破--过热--受挫--沉淀--再突破'' 的循环. 热潮会带来资源和想象力, 寒冬也会迫使研究者回到问题本身, 重新检查数据是否足够、模型是否合适、评价是否可靠、工程条件是否成熟. 因此, 学习机器学习不仅是在学习一组算法, 也是在学习如何把宏大的智能愿景落到可度量、可训练、可改进的系统上.
    \end{itemize}
\end{enumerate}

\textbf{问题思考}: 

\begin{itemize}
  \item 监督学习与无监督学习在训练数据形式和任务目标上有何根本区别?
  \item 回归任务与分类任务虽然都属于监督学习, 但它们在输出空间和评价方式上有何不同?
  \item 从两次 AI 寒冬看, 为什么仅靠人工规则难以支撑通用的智能系统?
\end{itemize}


\section{机器学习的四个基本要素}

\textbf{本节导读}:
\begin{itemize}
  \item \textbf{本节目的}: 把机器学习系统拆成数据、模型、学习准则和优化算法四个可分析的组成部分.
  \item \textbf{为什么学}: 后续每一种算法都可以看作这四个要素的不同选择; 先掌握这个框架, 才能比较不同方法到底改了哪里.
  \item \textbf{核心内容}: 数据提供经验来源, 模型规定可选函数集合, 学习准则定义什么叫“好”, 优化算法负责在模型集合中搜索参数.
  \item \textbf{最小例子}: 做房价预测时, 历史房屋记录是数据, 线性回归是模型, 均方误差是学习准则, 梯度下降或正规方程是寻找参数的方法.
  \item \textbf{最该记住}: 机器学习不是单独一个模型公式, 而是“数据 + 模型 + 损失 + 优化”的完整闭环.
\end{itemize}

\begin{enumerate}
  \item \textbf{数据}
    \begin{itemize}
      \item 数据是模型学习的基础, 其质量、规模与分布特性直接影响模型的最终性能.
      \item 训练集用于学习模型参数, 是模型拟合规律的主要依据.
      \item 验证集用于模型选择、超参数调整和训练过程监控, 其核心作用是评估模型的泛化能力.
      \item 测试集用于在训练完成后给出最终性能评估, 一般不参与模型设计和参数调节.
    \end{itemize}
  \item \textbf{模型}
    \begin{itemize}
      \item 模型是对输入与输出关系的数学刻画, 决定了学习系统能够表达何种规律.
      \item 线性模型假设输入与输出之间满足较简单的线性或广义线性关系, 例如线性回归和逻辑回归.
      \item 非线性模型通过特征变换、核方法或神经网络等手段增强表示能力, 从而适应更复杂的数据分布与决策边界.
    \end{itemize}
  \item \textbf{学习准则}
    \begin{itemize}
      \item 损失函数用于刻画单个样本上的预测误差, 常写为 $\mathcal{L}(y, f(x))$, 例如均方误差和交叉熵损失.
      \item 期望风险描述模型在真实数据分布上的平均损失:
      \[
        R(f) = \mathbb{E}[\mathcal{L}(y, f(x))].
      \]
      \item 由于真实分布通常未知, 实际训练中常采用训练集上的经验风险作为近似:
      \[
        \hat{R}_D(f) = \frac{1}{N} \sum_{i=1}^{N} \mathcal{L}(y_i, f(x_i)).
      \]
      \item 经验风险最小化原则通过最小化 $\hat{R}_D(f)$ 来寻找模型, 但若模型容量过强或数据不足, 仅追求经验风险最小可能导致过拟合.
    \end{itemize}
  \item \textbf{优化算法}
    \begin{itemize}
      \item 优化算法负责根据学习准则搜索合适的模型参数, 是连接模型与数据的计算机制.
      \item 梯度下降法利用损失函数关于参数的梯度信息, 通过迭代方式不断减小目标函数值.
      \item 随机梯度下降及其变体每次使用一个样本或一个小批量样本更新参数, 在大规模数据场景下更为常用.
      \item 学习率控制参数更新的步长; 学习率过大可能导致震荡甚至发散, 过小则会使收敛速度明显下降.
    \end{itemize}
\end{enumerate}

\textbf{问题思考}: 

\begin{itemize}
  \item 为什么模型评估不能只依赖训练集结果?
  \item 损失函数的定义会如何影响模型训练目标与最终行为?
  \item 与标准梯度下降相比, 随机梯度下降在计算效率和收敛性质上有哪些典型差异?
\end{itemize}


\section{泛化能力与正则化}

\textbf{本节导读}:
\begin{itemize}
  \item \textbf{本节目的}: 说明为什么模型最终要看未见数据上的表现, 并引入过拟合、欠拟合、验证集和正则化.
  \item \textbf{为什么学}: 真实任务中训练集只是样本, 如果只追求训练集误差最低, 模型可能学到噪声而不是规律.
  \item \textbf{核心内容}: 泛化能力衡量模型从训练数据推广到真实分布的能力, 正则化则通过限制模型复杂度或干预训练过程来改善这种能力.
  \item \textbf{最小例子}: 用高阶多项式拟合少量带噪声数据时, 曲线可以穿过所有训练点, 但在新点上剧烈摆动; 加入 $L_2$ 正则化或提前停止可以让曲线更平滑.
  \item \textbf{最该记住}: 训练误差低不等于模型好, 真正重要的是验证集和测试集上的稳定表现.
\end{itemize}

\begin{enumerate}
  \item \textbf{泛化与过拟合}
    \begin{itemize}
      \item 泛化能力是指模型在未见样本上的表现能力, 它是衡量机器学习系统是否有效的核心标准之一.
      \item 过拟合是指模型在训练集上表现很好, 但在新数据上性能下降的现象; 常见原因包括模型过于复杂、训练数据不足、数据噪声较大或训练时间过长.
      \item 与之相对, 欠拟合表示模型表达能力不足, 无法充分刻画数据中的主要规律.
    \end{itemize}
  \item \textbf{正则化的基本思想}
    \begin{itemize}
      \item 正则化的目标是在拟合训练数据与控制模型复杂度之间取得平衡, 从而提升模型对未知数据的泛化能力.
      \item 从模型复杂度控制的角度看:
        \begin{itemize}
          \item $L_1$ 正则化倾向于得到稀疏参数表示, 因而常用于特征选择或稀疏建模.
          \item $L_2$ 正则化倾向于抑制参数过大, 有助于提升模型稳定性与泛化性能.
          \item 数据增强通过构造更多合理训练样本来扩展数据分布覆盖范围, 常用于图像、语音和文本任务.
        \end{itemize}
      \item 从训练过程干预的角度看:
        \begin{itemize}
          \item 提前停止利用验证集性能决定训练终止时机, 避免模型继续拟合训练集中的偶然噪声.
          \item 权重衰减通常可看作 $L_2$ 正则化在优化过程中的具体实现方式之一.
          \item Dropout 通过在训练阶段随机屏蔽部分神经元或连接, 降低网络对局部共适应的依赖.
        \end{itemize}
    \end{itemize}
  \item \textbf{验证集的作用}
    \begin{itemize}
      \item 验证集用于比较不同模型或超参数配置在未参与训练的数据上的表现, 是模型选择的重要依据.
      \item 在实际训练中, 验证集还常用于监控过拟合趋势, 例如决定是否提前停止训练.
      \item 因此, 验证集并不是用于最终报告性能, 而是服务于训练过程中的选择与调优.
    \end{itemize}
\end{enumerate}

\textbf{问题思考}: 

\begin{itemize}
  \item 过拟合通常由哪些因素共同导致? 这些因素之间是否会相互影响?
  \item $L_1$ 与 $L_2$ 正则化分别更适合解决哪些问题?
  \item 为什么验证集能够帮助我们提升模型的泛化能力, 但测试集不能被频繁用于调参?
\end{itemize}





(注: 详细定义和更多的内容请参考课程教材第一、二章)

\end{document}
